\section{An actual-count lower bound for the small-index schedule}
\label{subsec:obstruction}

The following result concerns exactly the interpolation count and catalyst
decomposition in \eqref{eq:relative-counts}. Its hypothesis uses the
actual branching count, rather than the sufficient upper bound
\eqref{eq:finite-exponent}.
We use \(\log_4 48\), the recursive exponent of a known rational
rank-48 algorithm for \(4\times4\) multiplication~\citep{Dumas26},
as a fixed comparison value.

\begin{proposition}[A limitation of parameter tuning]\label{prop:obstruction}
For the \(a=2\) construction, use \(E=kD+1\), the catalyst bound
\eqref{eq:small-kappa}, and the full-batch recursion of
Lemma~\ref{lem:batch}. Assume \(B>A\). If
\[
 \frac{\log\lceil r/m\rceil}{\log(B/A)}
 <\rho_{48}:=\frac{\log48}{\log4},
\]
then its recursive input dimension \(A\) has more than
\(11.37\) billion decimal digits. One of its separator sector counts
also exceeds \(10^{1206}\).
\end{proposition}
\begin{proof}
Use the small-index counts above and define
\begin{align*}
 \delta&=30\rho_{48}\log(10/3)-90\log3,\\
 \gamma&=6(35\log3+56\log4+54\log5+44\log6+26\log7+15\log8),\\
 \mu&=\log\binom{56}{36}.
\end{align*}
Thus \(\log\kappa_N=N\gamma\).
The hypothesis and \(\lceil r/m\rceil\ge r/m\) imply the necessary
condition
\begin{equation}\label{eq:necessary-cost}
 \frac{\log E}{k}+\frac{N\gamma}{j}
 <\rho_{48}\Delta_N-90N\log3.
\end{equation}
We have dropped the positive copy-ratio charge, making the condition weaker.

The one-sided Stirling estimate~\citep{Robbins55,DLMF26}
\(\log(n!)=(n+\tfrac12)\log n-n+\tfrac12\log(2\pi)+r_n\),
\(0<r_n<1/(12n)\), gives
\begin{equation}\label{eq:entropy-loss-upper}
 \Delta_N\le30N\log(10/3)-7\log N-2C_0+\frac{2\eta}{N},
\end{equation}
where
\begin{align*}
 C_0&=\frac12\sum_{b=2}^6\log(6u_bd_b/w_b)
       +\log6+\frac32\log5-\frac12\log15,\\
 \eta&=\sum_{b=2}^6\frac1{12w_b}+\frac1{180}.
\end{align*}
For each binomial coefficient, subtract the three Stirling formulas and
keep only the positive remainder \(1/(12Nw_b)\). Its correction to
entropy is at most
\(-\frac12\log N-\frac12\log(2\pi u_bd_b/w_b)+1/(12Nw_b)\).
The three-sector coefficient contributes
\(-\log N-\log(2\pi)+\frac12\log15-\frac32\log5+1/(180N)\).
Replace \(2\pi\) by the smaller number six and use the exact entropy
identity to obtain \eqref{eq:entropy-loss-upper}.

The right side of \eqref{eq:necessary-cost} is therefore at most
\[
 f(N)=\delta N-\rho_{48}(7\log N+2C_0-2\eta/N).
\]
Its second derivative is
\(\rho_{48}(7/N^2+4\eta/N^3)>0\).
The archived rational enclosures from \eqref{eq:log-enclosure} give
\(f(1)<-71.4589\) and \(f(80)<-0.224635\).
A convex function is at most its endpoint chord on \([1,80]\).
The necessary condition has a positive left side, so \(N\ge81\).
The penalty \(7\log N+2C_0-2\eta/N\) is increasing. Put
\(\beta=\rho_{48}(7\log81+2C_0-2\eta/81)\). Then
\begin{equation}\label{eq:necessary-simplified}
 \frac{\log E}{k}+\frac{N\gamma}{j}<\delta N-\beta.
\end{equation}

Selecting 36 positions within each of \(N\) disjoint blocks of 56 proves
\(M_3=\binom{56N}{36N}\ge\binom{56}{36}^N\).
Because \(M_3\ge2\),
\(E\ge6(M_3-1)^2+1\ge M_3^2\), hence \(\log E\ge2N\mu\).
Arithmetic--geometric mean in \eqref{eq:necessary-simplified} yields
\[
 Nkj>\frac{8\mu\gamma N^3}{(\delta N-\beta)^2}
       \ge\frac{54\mu\gamma\beta}{\delta^3}.
\]
Here \(\delta N>\beta\), and the last inequality follows from
\(4x^3-27(x-1)^2=(x-3)^2(4x-3)\ge0\), \(x=\delta N/\beta>1\).
The same saved rational enclosures certify
\[
 \frac{54\mu\gamma\beta}{\delta^3}>82\,111\,374.98,
 \qquad
 460\frac{54\mu\gamma\beta}{\delta^3}\frac{\log2}{\log10}
       >11\,370\,273\,954.02,
 \qquad 81\frac\mu{\log10}>1206.
\]
Since \(A=2^{460Nkj}\), these inequalities prove both assertions.
\end{proof}

The conclusion does not apply after pruning zero branches, replacing the
catalyst decomposition, sharing bilinear work, or changing exactification.
Those operations change the counted schedule. Factorized Fourier circuits
can reduce its linear work and storage, but by themselves do not reduce the
independent source calls counted in \eqref{eq:relative-counts}.
